% !TEX program = xelatex
% 张宇考研数学预测8套卷（2027年，数学一，卷一），原试卷第5～8页。
\documentclass[zihao=5,fontset=windows,oneside]{exam-zh}
\examsetup{
  page/size=a4paper,
  page/show-head=false,
  page/foot-content={第;页（共;页）},
  sealline/show=false,
  question={hang=true,show-answer=false,top-sep=5pt,bottom-sep=3pt},
  choices={columns=0,column-sep=1.6em,linesep=5pt,label-pos=top-left},
  fillin={no-answer-type=none,width=19mm},
}
\linespread{1.1}
\setCJKfamilyfont{examcoverfont}[BoldFont=FandolSong-Bold.otf]{FandolSong-Regular.otf}
\providecommand{\dd}{\,\mathrm d}
\newcommand{\WorksheetTitle}[1]{%
  \begin{center}
    {\CJKfamily{examcoverfont}\bfseries\fontsize{17}{22}\selectfont 2027年张宇考研数学预测8套卷}\\[5pt]
    {\CJKfamily{examcoverfont}\fontsize{13}{18}\selectfont 数学一 · 第1套 · #1}
  \end{center}\vspace{8pt}%
}
\newcommand{\QuestionBreak}{}
\ifdefined\QuestionCrops
  \geometry{textwidth=468bp,hcentering}
  \renewcommand{\QuestionBreak}{\clearpage}
  \AtBeginDocument{\pagestyle{empty}}
\fi

\examsetup{question={show-points=true,points-prelabel={（本题满分},points-postlabel={分）},points-separate-par=true}}
\AtBeginDocument{\hypersetup{pdftitle={2027年张宇考研数学预测8套卷 · 数学一 · 第1套 · 解答题}}}
\newcommand{\WorkSpace}{\ifdefined\QuestionCrops\else\vfill\fi}
\newcommand{\WorkPageBreak}{\ifdefined\QuestionCrops\else\clearpage\fi}
\begin{document}

\ifdefined\QuestionCrops\else
  \WorksheetTitle{解答题}
  \section{解答题：第17～22题，共70分。解答应写出文字说明、证明过程或演算步骤。}
\fi

\begin{question}[index=17,points=10]
求函数 $f(x,y)=(y-x)(y-x^2)$ 的极值。
\end{question}
\QuestionBreak\WorkSpace

\begin{question}[index=18,points=12]
已知函数 $g(x)$ 连续，$g(0)=0$，设 $\displaystyle f(x)=\int_0^x g(x^2t)\dd t$，求 $f'(x)$ 的表达式，并判断 $f'(x)$ 在 $x=0$ 处的连续性。
\end{question}
\QuestionBreak\WorkSpace\WorkPageBreak

\begin{question}[index=19,points=12]
若曲线 $y=f(x)$ 的切线斜率在 $[a,b]$ 上单调递增，证明曲线 $y=f(x)$ 在 $[a,b]$ 上是凹的。
\end{question}
\QuestionBreak\WorkSpace

\begin{question}[index=20,points=12]
若对第一卦限内任一有向闭曲面 $\Sigma$ 恒有
\[
  \oiint_{\Sigma} x^2z\varphi(y)\dd y\dd z
  +2xz\varphi'(y)\dd z\dd x-xz^2\mathrm e^y\dd x\dd y=0,
\]
其中 $\varphi(y)$ 具有二阶连续导数，
$\displaystyle\varphi(\pi)=\varphi'(\pi)=\frac{\mathrm e^\pi-1}{2}$，求 $\varphi(y)$。
\end{question}
\QuestionBreak\WorkSpace\WorkPageBreak

\begin{question}[index=21,points=12]
设 $\displaystyle
  \symbfit\alpha_1=\begin{pmatrix}2\\2\\-1\end{pmatrix},\quad
  \symbfit\alpha_2=\begin{pmatrix}-1\\-3\\2\end{pmatrix},\quad
  \symbfit\alpha_3=\begin{pmatrix}0\\-4\\3\end{pmatrix}$，并记
$\displaystyle\symbfit A=\begin{pmatrix}2&-1&0\\2&-3&-4\\-1&2&3\end{pmatrix}$。
\begin{enumerate}[label=（\arabic*）,leftmargin=2.5em,labelsep=.4em,itemsep=3pt,topsep=3pt]
\item 求 $\symbfit\alpha_1,\symbfit\alpha_2,\symbfit\alpha_3$ 的一个极大线性无关组；
\item 求 $\symbfit A^n\ (n>2)$。
\end{enumerate}
\end{question}
\QuestionBreak\WorkSpace

\begin{question}[index=22,points=12]
设随机变量 $X$ 的概率分布为 $\displaystyle P\{X=k\}=\frac{1}{2^{k+1}}+\frac{a}{3^k}$，$k=1,2,\dots$。
\begin{enumerate}[label=（\arabic*）,leftmargin=2.5em,labelsep=.4em,itemsep=3pt,topsep=3pt]
\item 求 $a$ 的值；
\item 求 $X$ 的数学期望 $EX$ 与方差 $DX$；
\item 比较 $E(X\mid X>1)$ 与 $E(X\mid X>2)$ 的大小，并说明理由。
\end{enumerate}
\end{question}
\WorkSpace

\end{document}
